(1) 对理想气体等温自由体积扩张，$dS=Nk_BdV/V$。A 气体从 $V_A$ 扩张至 $V$、B 气体从 $V_B$ 扩张至 $V$，由不同物种的熵可加得 $\Delta S=k_B[N_A\ln(V/V_A)+N_B\ln(V/V_B)]>0$。若初态每组都处于相同压强，$V_A/V=N_A/(N_A+N_B)$，可改写为标准 $-k_B(N_A\ln x_A+N_B\ln x_B)$。
(2) 混合后各组分在整个体积均匀分布，因此分压 $p_A=N_Ak_BT/V$、$p_B=N_Bk_BT/V$，总压 $p=(N_A+N_B)k_BT/V$。理想气体化学势在定温下含 $\mu_i(T,p_i)=\mu_i^\circ(T)+k_BT\ln(p_i/p^\circ)$，故组分 $i$ 的化学势变化为 $k_BT\ln(V_i/V)<0$；和粒子数相乘得到总 Gibbs 自由能下降 $\Delta G=-T\Delta S$。
(3) 混合过程与可逆分离过程之间的可用最大等温功为 $W_{\rm out,max}=-\Delta G=T\Delta S=k_BT[N_A\ln(V/V_A)+N_B\ln(V/V_B)]$，这里要求理想选择透过膜，且混合前后外界维持相同温度。直接去隔板的不可逆自由混合不作功，因而损失了这部分自由能。
(4) 若 A、B 实为完全相同且不可辨别的粒子，去隔板前后宏观平衡态没有不同物种浓度差，不能用有标签的组分自由扩张量重复计数微观态。用正确的不可区分粒子计数，在初始温度和压强相同的前提下无真正混合熵：$\Delta S_{\rm mix}=0$，可提取的理想选择性混合功也为零。若两侧初态压强不同，仍可能有压力平衡熵产，不能混同于物种混合熵。
